% s-176-p176-m003; original PDF 176, 177
\Aphorism{s-176-p176-m003}{265}
Ich verstehe es vollkommen, wenn der Europäer in unseren Breiten eine neue Krankheit, die \Sperr{Wetter-Neurasthenie}, bekommt: er reagirt damit nur auf das immer entschiedener hervortretende Missverhältniss zwischen Klima und Organismus und empfängt als haruspicium (ach, aus seinen eigenen Eingeweiden!) die dringende Mahnung, nach milderen Erdstrichen auszuwandern. Unser Frühling kündigt sich durch eine Sturmfluth von Todesanzeigen in den Zeitungen an; diese Art „Lenz“ versteht es mit ihren Vorsommertagen und Winterrückschlägen, den Leidenden, die sich schon wieder für ein Jahr gefristet wähnen; auf heimtückische Weise den Rest zu geben. Folgt der Sommer, der zwischen dumpfer Hitze und Regen derart abzuwechseln pflegt, dass der arbeitende Theil der Bevölkerung um seine Erholung kommt: Abende, Sonntage und Ferienwochen sind unwiderruflich nass, wüst und unerquicklich. Der Herbst übernimmt die Regenerbschaft des Sommers und reicht damit bis December oder Januar aus; dann ein Spätwinter ohne Charakter, wenig Schnee, aber viel kalter Schlamm, und wieder Frühling. Das Ganze ein System der Systemlosigkeit, ein wohldurchdachter Unsinn, ein Zufallsspiel mit ausschliesslichen Verlustchancen, ein \Foreign{modus vivendi}, bei dem Niemand leben kann. Dem Landwirth wächst das Korn aus, dem voraus disponirenden Kaufmann wird seine Disposition zu Schanden, dem Berggastwirth verregnen seine drei Monate Saison — und der Eine wie der Andere trägt sein Maximum an leichten und schweren Erkältungen davon. Das nennt man Anpassung der Organismen an die klimatischen Bedingungen, das Ergebniss etlicher Jahrtausende biologischer Zuchtwahl. Thatsache ist, dass wir gar nicht angepasst sind, dass wir in einem unbewohnbaren Erdtheil Fuss gefasst haben und gedankenloser Weise wohnen bleiben, statt ihn den Grönländern zu überlassen; diesen ungeheuren Missgriff bemänteln wir mit der wissenschaftlichen Illusion, uns nur in einer vorübergehenden Nässeperiode zu befinden und wieder wärmeren, trockneren Jahreszeiten entgegenzugehen. Thorheit! Mitteleuropa ist die Urheimat der Nässe, der ewige Regenstrich und Port Tarascon der nördlichen Erdhalbkugel; das war nie anders und wird nicht anders, und wenn der Mitteleuropäer nicht auswandert, hat er die Wahl, an der Nässe zu Grunde zu gehen oder am Gegenmittel gegen die Nässe, am Alkohol!
\par

